%==============================================================================
% Overleaf-ready LaTeX file — REORGANISED into a thesis-style chapter/section
% hierarchy. Preserves your existing prose but restructures it.
%
% Where a section can be filled verbatim from your existing
% 09_overleaf_ready.tex, you will see one of two markers:
%
%   * "% INSERT_FROM_ORIGINAL: lines X--Y" — copy that block from the
%     original file verbatim into this position.
%   * "% NEW_CONTENT_REQUIRED:" — section that is new relative to the
%     original; a one-paragraph stub is provided that you should expand.
%
% Top-level structure (mapped to user's content):
%   1. Introduction              — § Introduction + Background + Related Work + Dataset § (condensed)
%      1.1 DermaMNIST Dataset
%      1.2 Pipeline Design and Overview
%   2. Method                    — § Experimental Setting (renamed + restructured)
%      2.1 Dataset Preprocessing and Augmentation
%      2.2 Centralised Training Framework
%      2.3 Training and Evaluation Protocols
%      2.4 Federated Learning Framework
%          2.4.1 Baseline Federated Setup
%          2.4.2 Heterogeneity Scenarios in Federated Learning
%          2.4.3 FedProx
%          2.4.4 FedNova
%   3. Results and discussion    — § Results + § Discussion (interleaved)
%      3.1 Centralised baseline
%      3.2 Federated Learning
%          3.2.1 Initial Setup: Engineered Paired Partition (n=10 seeds)
%          3.2.2 Statistical Heterogeneity (multi-partition: IID, Dirichlet, specialist, paired)
%          3.2.3 System Heterogeneity (C0, C1, C2)
%          3.2.4 Summary and Comparison of Heterogeneity Experiments
%      3.3 Hardware and Resource Limitations
%   4. Conclusion                — § Conclusion + § Future Work folded in
%   References
%==============================================================================
\documentclass[11pt,a4paper]{report}  % NB: 'report' supports \chapter; was 'article'

\usepackage[utf8]{inputenc}
\usepackage[T1]{fontenc}
\usepackage{graphicx}
\usepackage{booktabs}
\usepackage{amsmath, amssymb}
\usepackage{microtype}
\usepackage{hyperref}
\usepackage{caption}
\usepackage{xcolor}
\usepackage[margin=2.5cm]{geometry}
\usepackage{setspace}
\usepackage{tikz}
\usetikzlibrary{positioning, shapes.geometric, arrows.meta, fit, backgrounds}
\usepackage{algorithm}
\usepackage{algpseudocode}
\usepackage[authoryear,round]{natbib}
\onehalfspacing

% --- TODO / pending-HPC markers ----------------------------------------------
\newif\ifshowtodo
\showtodotrue
\newcommand{\TODOhpc}[1]{%
  \ifshowtodo\textcolor{red}{\textbf{[HPC PENDING: #1]}}\else#1\fi}
\newcommand{\PENDING}{\TODOhpc{result}}

\title{Federated Optimisation on DermaMNIST: A Paired-Seed Empirical
Comparison of FedAvg, FedProx, and FedNova Under Statistical and
Computational-Work Heterogeneity}
\author{Barbara Koch}
\date{}

\begin{document}

%==============================================================================
% TITLE PAGE
%==============================================================================
\begin{titlepage}
\centering
\vspace*{3cm}
{\Huge\bfseries Federated Optimisation on DermaMNIST\par}
\vspace{0.5cm}
{\Large A Paired-Seed Empirical Comparison of FedAvg, FedProx, and FedNova\par}
{\Large Under Statistical and Computational-Work Heterogeneity\par}
\vspace{4cm}
% \includegraphics[width=4cm]{cambridge_crest.png}  % uncomment if you have the crest
\vspace{2cm}
{\large\bfseries Barbara Koch\par}
\vspace{1cm}
{\large Supervisor: \TODOhpc{supervisor name}\par}
{\large MPhil in Data Intensive Science\par}
{\large University of Cambridge\par}
\vspace{3cm}
{This dissertation is submitted for the degree of\par}
{\itshape Master of Philosophy\par}
\vfill
{\large \TODOhpc{College name} \hfill \TODOhpc{Month} 2026 \par}
\end{titlepage}

%==============================================================================
% DECLARATION
%==============================================================================
\chapter*{Declaration}
\addcontentsline{toc}{chapter}{Declaration}
\thispagestyle{empty}

I hereby declare that, except where specific reference is made to the work of
others, the contents of this dissertation are original and have not been
submitted in whole or in part for consideration for any other degree or
qualification in this, or any other university. This dissertation is my own
work and contains nothing which is the outcome of work done in collaboration
with others, except as specified in the document. During its preparation,
GitHub Copilot and Claude Code were utilised to assist with code organisation,
autocompletion, experimental orchestration, and verifying results. For
language accuracy, Grammarly was employed to check spelling and grammar
throughout the document. This dissertation contains \TODOhpc{word count}
words, excluding references and declaration.

\vspace{2cm}
\hfill Barbara Koch

\hfill \TODOhpc{Month} 2026

%==============================================================================
% ABSTRACT
%==============================================================================
\chapter*{Abstract}
\addcontentsline{toc}{chapter}{Abstract}
\thispagestyle{empty}

Federated Learning (FL) trains models across data silos without
centralising patient data, but the canonical FedAvg algorithm
\citep{mcmahan2017fedavg} suffers under client heterogeneity. We
evaluate FedProx \citep{li2020fedprox}, which adds a proximal
regulariser to local objectives, on DermaMNIST --- a 7-class
skin-lesion classification dataset with severe class imbalance
(58:1 between majority and minority classes). We design a custom
non-IID partition (\texttt{balanced\_paired\_7\_clients}) in which
every minority class is held by exactly two clients, avoiding the
catastrophic per-class regression observed under singleton-class
partitions. Across 10 paired-seed runs ($R = 150$ rounds, $E = 20$
local epochs, $\mu = 0.01$), FedProx achieves a mean test macro-F1
of $0.508 \pm 0.014$ against FedAvg's $0.481 \pm 0.025$ --- a
within-pair improvement of $\Delta = +0.027 \pm 0.035$ (paired
Wilcoxon $p = 0.020$, rank-biserial $r = +0.818$). FedProx wins on
9 of 10 paired seeds. The largest per-class gain is on melanoma
($\Delta = +0.114$, $p = 0.006$), the most clinically critical
lesion. No class regresses by more than $0.012$ F1, satisfying the
pre-registered safety criterion. FedProx additionally reduces
across-seed variance by 45\%, suggesting more stable convergence
under non-IID conditions. A third algorithm, FedNova
\citep{wang2020fednova}, is implemented and evaluated under
computational-work heterogeneity as the canonical comparator for
heterogeneous local SGD step counts.

\vspace{1em}
\noindent\textbf{Keywords:} federated learning, FedAvg, FedProx, FedNova,
DermaMNIST, dermatology, class imbalance, non-IID, paired-seed protocol,
medical imaging, data privacy

%==============================================================================
% TOC / LoF / LoT
%==============================================================================
\tableofcontents
\listoffigures
\listoftables

%==============================================================================
\chapter{Introduction}
\label{ch:intro}
%==============================================================================

% ============================================================================
% INTRO CHAPTER STRUCTURE:
% - Three-to-four opening paragraphs on the clinical and methodological
%   motivation (skin-cancer + dermatoscopic imaging + privacy + FL).
% - Brief overview of the FedAvg → non-IID problem → FedProx/FedNova
%   remedies (condensed from your current §Background §sec:background
%   and §Related Work §sec:related — lit review folded INTO the
%   Introduction rather than given a separate chapter).
% - §1.1 Dataset — short standalone description.
% - §1.2 Pipeline Design and Overview — short, with the pipeline figure.
% ============================================================================

% INSERT_FROM_ORIGINAL: lines 100-217 of 09_overleaf_ready.tex
% (Your existing opening narrative: skin cancer + HAM10000 + GDPR/HIPAA +
%  FedAvg/non-IID + FedProx/FedNova + scope conditions. Keep it almost
%  verbatim, just remove the explicit \section{Introduction} header
%  since this is now a \chapter.)

Skin cancer is one of the most common malignancies worldwide, and
melanoma in particular accounts for the great majority of skin-cancer
deaths despite representing fewer than 5\% of cases. The outcome of a
melanoma diagnosis is exquisitely sensitive to the time at which the
disease is detected: five-year survival rates exceed 99\% when the
lesion is identified while still localised, falling to under 35\% once
the disease has metastasised. This stark survival gradient has
motivated a generation of computer-aided diagnostic systems based on
dermoscopic images, and the recent availability of large-scale
labelled image datasets has made deep learning the de facto
methodology in this space. The HAM10000 collection
\citep{tschandl2018ham10000}, assembled from multiple clinical sites
and comprising 10{,}015 dermoscopic images across seven diagnostic
categories, has become the standard benchmark for this class of
system.

% INSERT_FROM_ORIGINAL: lines 113-217 (motivation continuation through
% to scope conditions). Cut at "...identified as the primary direction
% for follow-up work in §\ref{sec:future}." and pick up below.

% ----- Condensed literature review (folded into Intro) -----
% The original draft has two separate chapters (§Background and §Related
% Work) totalling ~525 lines. Recommended condensation:
%
%   (a) One paragraph on FedAvg + the non-IID degradation problem
%       (condense from lines 553-579).
%   (b) One paragraph on the three remedy families (proximal / control-
%       variates / aggregation-normalisation) — condense lines 581-631.
%   (c) One paragraph on non-IID benchmarks (Dirichlet, NIID-Bench) —
%       condense lines 660-727.
%   (d) One paragraph on medical-imaging FL (Sheller, Rieke, FeTS, FedBN)
%       — condense lines 732-822.
%   (e) One paragraph stating where this thesis fits (lines 804-822).
%
% Each paragraph 4-7 sentences. Total ~1.5 pages.

% NEW_CONTENT_REQUIRED: at the end of the Introduction, before §1.1,
% drop in your "Contributions" itemised list from lines 219-275 of the
% original.

\paragraph{Contributions.} The substantive contributions of this
dissertation are:
\begin{itemize}
\item \textbf{Paired-seed protocol with bit-equivalence at $\mu = 0$.}
We design an experimental protocol in which FedAvg and FedProx share
identical partitions, initial model weights, per-(round, client)
DataLoader random-number generators, dropout masks, and
system-heterogeneity schedules. Bit-equivalence at $\mu = 0$ is
verified by seven unit tests at $10^{-8}$ tolerance over per-round
history, final aggregated state-dictionaries, and test-set metrics.

\item \textbf{A custom \texttt{balanced\_paired\_7\_clients} partition
supporting a clean per-class mechanism prediction.} The partition is
designed so that every minority class is held by exactly two of seven
clients while the majority class (melanocytic nevi) is held by all
seven.

\item \textbf{Headline empirical result with effect-size and
leave-one-seed-out robustness.} Across ten paired seeds, FedProx
($\mu = 0.01$) improves over FedAvg on test macro-F1 by
$\Delta = +0.027 \pm 0.035$ (paired Wilcoxon $p = 0.020$,
sign-test $p = 0.021$, rank-biserial $r = +0.818$, 9 of 10 paired
wins).

\item \textbf{Three-regime robustness analysis.} The headline finding
is contextualised by sweeps under three statistical-heterogeneity
regimes: the custom paired partition, a uniform IID null-falsification
control, and the literature-standard Dirichlet-$\alpha = 0.1$
non-IID partition.

\item \textbf{Independently verified FedNova implementation.} FedNova
\citep{wang2020fednova} is verified against canonical reference values
(e.g.\ $a_i = 41.381$ for $\tau = 10$, $m = 0.9$) and used in a
three-way paired-seed comparison on DermaMNIST.
\end{itemize}

%==============================================================================
\section{DermaMNIST Dataset}
\label{sec:dataset-intro}
%==============================================================================

% INSERT_FROM_ORIGINAL: lines 906-935 (Dataset section from original) +
% lines 1271-1310 (the "simulated, not real" methodology framing).

We use DermaMNIST \citep{yang2023medmnistv2}, the dermatology subset of
the MedMNIST~v2 benchmark, derived from the HAM10000 collection of
dermatoscopic images \citep{tschandl2018ham10000}. DermaMNIST contains
10{,}015 dermoscopic images of skin lesions across seven diagnostic
categories, split into 7{,}007 training, 1{,}003 validation, and 2{,}005
test samples in the canonical MedMNIST split. The source archive
\texttt{dermamnist\_64.npz} (the v2 distribution at 64$\times$64 RGB) is
the canonical MedMNIST file format; we resize all images to
28$\times$28 before training. The class distribution is severely
imbalanced, with a 58:1 ratio between the most prevalent
(\textit{melanocytic nevi}, 67.05\%) and least prevalent
(\textit{dermatofibroma}, 1.15\%) classes. This imbalance is preserved
without re-sampling so that any reported gain reflects realistic
clinical priors.

% INSERT_FROM_ORIGINAL: lines 1271-1310 here — the "Simulated, not
% real, heterogeneity" scope paragraph and the bulleted list of
% heterogeneity axes not modelled.

%==============================================================================
\section{Pipeline Design and Overview}
\label{sec:pipeline}
%==============================================================================

% INSERT_FROM_ORIGINAL: lines 839-904 (the pipeline TikZ figure).

In summary, this dissertation presents a controlled empirical pipeline
for paired-seed federated optimisation on a non-IID classification
benchmark. First, we establish a centralised baseline on the pooled
DermaMNIST training set to anchor the upper bound on global-model
performance. Second, we implement a 7-client federated-learning
framework under a custom \texttt{balanced\_paired} partition designed
to maximise the proximal-anchor mechanism's engagement. Third, we
investigate the FedProx-over-FedAvg gap under three forms of
statistical heterogeneity (IID, Dirichlet-$\alpha = 0.1$, and the
custom paired partition) and three forms of computational-work
heterogeneity (uniform, fixed stragglers, random stragglers). Finally,
we introduce FedNova \citep{wang2020fednova} as the canonical
comparator for heterogeneous local SGD step counts. This work aims to
provide a paired-seed inferential answer to the FedAvg-vs-FedProx
question on DermaMNIST, addressing a gap left open by NIID-Bench
\citep{li2022niidbench}.

% INSERT_FROM_ORIGINAL: lines 839-904 (the full \begin{figure}...
% \end{figure} block containing the TikZ pipeline diagram).

%==============================================================================
\chapter{Method}
\label{ch:method}
%==============================================================================

This section details the implementation of the centralised baseline
and federated-learning frameworks. It encompasses dataset
preprocessing, model architecture, training and evaluation protocols,
and the three federated optimisation algorithms.

%==============================================================================
\section{Dataset Preprocessing and Augmentation}
\label{sec:dataset-preproc}
%==============================================================================

% NEW_CONTENT_REQUIRED: short standalone preprocessing section
% (about half a page).

All images and corresponding labels were taken directly from the
DermaMNIST v2 distribution \citep{yang2023medmnistv2}. To facilitate
faster experimentation and to align with prior FL-on-DermaMNIST
literature, all images were downsampled from $64 \times 64$ to
$28 \times 28$ pixels using bilinear interpolation. Pixel intensities
were normalised to the range $[0, 1]$. No data augmentation was
applied during federated training: the deliberate design choice is to
isolate the algorithmic comparison between FedAvg, FedProx, and
FedNova from the confound of stochastic augmentation, which would
otherwise interact with the per-(round, client) random-number-
generator seeding protocol described in §\ref{sec:protocol}. The
centralised baseline (§\ref{sec:centralised-results}) was likewise run
without augmentation so that the federation tax estimate is internally
consistent.

For all experiments, images are loaded in batches of size $B = 32$,
capped per-client at $\min(32, n_i)$ so that the smallest client
(C6 with $n = 673$ samples) still produces meaningful mini-batches.

%==============================================================================
\section{Centralised Training Framework}
\label{sec:centralised-method}
%==============================================================================

% NEW_CONTENT_REQUIRED: short paragraph describing the centralised
% baseline architecture and training protocol.

The classifier is a four-block convolutional neural network
(\texttt{DermMNISTCNN}, 423{,}175 parameters) with feature widths
$32 \to 64 \to 128 \to 256$. Each block consists of a single
3$\times$3 convolution (\texttt{padding=1}) followed by GroupNorm,
ReLU, and 2$\times$2 max-pooling; the final block replaces max-pooling
with adaptive global average pooling. A fully connected head of
$256 \to 128 \to 7$ with dropout~0.2 produces the class logits.
% INSERT_FROM_ORIGINAL: lines 944-959 (the verbatim block describing
% the CNN graph + the GroupNorm justification citing
% \citet{wu2018groupnorm} and \citet{hsieh2020quagmire}).

The centralised baseline is computed by training this same
\texttt{DermMNISTCNN} architecture on the pooled 7{,}007-sample training
set under the same evaluation protocol used for the federated runs
(early stopping on validation macro-F1; test evaluated once at the
best-val epoch). The centralised model is trained for 50 epochs with
SGD (learning rate 0.01, momentum 0.9, weight decay 0.0, batch size
32) and provides the upper bound on test macro-F1 achievable when no
client partition is imposed.

%==============================================================================
\section{Training and Evaluation Protocols}
\label{sec:training-protocols}
%==============================================================================

% NEW_CONTENT_REQUIRED + INSERT_FROM_ORIGINAL:
% (a) Brief paragraph describing optimiser/loss/eval.
% (b) Then insert from original lines 1387-1410 (Metrics and Statistical
%     Analysis).

The optimiser is mini-batch SGD with momentum coefficient $m = 0.9$
and no weight decay (set to 0 to avoid conflation with FedProx's
proximal regulariser). The loss function is the multi-class
cross-entropy, equivalent in form for all three algorithms with the
sole exception of FedProx's added proximal penalty
(§\ref{sec:fedprox}). Federated rounds use full participation
($C = 1.0$); $R = 150$ communication rounds are run per experiment at
$E = 20$ local epochs per round.

The primary outcome is test macro-F1, the unweighted mean of per-class
F1 scores. Accuracy is reported as a secondary descriptor but is
misleading under the severe class imbalance present in DermaMNIST.
Per-class F1 is reported separately to detect class-specific
regressions. Validation metrics are recorded after every communication
round and aggregated across the 10 paired seeds with mean and standard
error of the mean (SEM); the test set is evaluated once per run at the
round of best validation macro-F1, producing 20 single-shot test
evaluations from which all reported test statistics are computed.

% INSERT_FROM_ORIGINAL: lines 1398-1408 (paired Wilcoxon /
% rank-biserial / bootstrap-CI / per-class Holm description).

\subsection*{Paired-Seed Protocol}

% INSERT_FROM_ORIGINAL: lines 1354-1386 (the paired-seed protocol
% and reproducibility-provenance paragraph). This is a methodological
% device, not an experimental condition.

%==============================================================================
\section{Federated Learning Framework}
\label{sec:fl-framework}
%==============================================================================

% INSERT_FROM_ORIGINAL: lines 961-1078 (FL setup paragraphs + the
% hyperparameters table). Recommend keeping the FedAvg pseudocode block
% in this introductory text (before §2.4.1), then splitting FedProx /
% FedNova into their own subsections below.

The federated-learning protocol used throughout this dissertation
follows the canonical cross-silo formulation surveyed in
\citet{kairouz2021advances}. The headline statistical-heterogeneity
results were produced by the project's pure-PyTorch reference loop
(\texttt{experiments/run\_one.py}); a Flower~1.29 simulation runtime
\citep{beutel2020flower} via \texttt{flwr.simulation.start\_simulation}
(each client a Ray actor) is also implemented and is used for the
system-heterogeneity, IID, and Dirichlet sweeps. Cross-runtime
equivalence at full experimental scale is verified at two extreme
seeds (789 and 8675309) and reported as a sensitivity check.

% INSERT_FROM_ORIGINAL: rest of lines 961-1040 (cross-runtime
% framing, the hyperparameters table at lines 1054-1078, and the
% FedAvg aggregation equation around line 980).

%==============================================================================
\subsection{Baseline Federated Setup}
\label{sec:baseline-fl}
%==============================================================================

% INSERT_FROM_ORIGINAL: lines 1042-1078 (the protocol paragraphs:
% "The server orchestrates a fixed pool of seven clients…", through to
% the hyperparameters table). Plus the canonical FedAvg description from
% lines 1226-1228.

The baseline federated setup uses 7 clients participating fully each
round, $R = 150$ communication rounds, and $E = 20$ local epochs per
client per round. The aggregation rule is the canonical size-weighted
mean,
\begin{equation}
w^{(t+1)} \;=\; \sum_{k=1}^{K} \frac{n_k}{n}\, w_k^{(t+1)},
\label{eq:fedavg-method}
\end{equation}
where $K = 7$ is the number of clients, $n_k$ is the number of
training samples at client $k$, and $n = \sum_k n_k$. This is the
FedAvg rule of \citet{mcmahan2017fedavg} and is used unchanged by both
FedAvg and FedProx; FedNova replaces it with the normalised aggregator
in §\ref{sec:fednova}.

%==============================================================================
\subsection{Heterogeneity Scenarios in Federated Learning}
\label{sec:het-scenarios}
%==============================================================================

% Bold subheaders inline for the two heterogeneity types (Statistical
% and System), with brief paragraphs describing each scenario. No
% \subsubsection numbering — just bold inline subheaders.

\paragraph{Statistical heterogeneity (label-distribution skew).}

To evaluate the effect of label-distribution skew across clients,
this dissertation tests four partitions of the DermaMNIST training
set:

\begin{itemize}
\item \textbf{IID (falsification check).} Each client receives a
uniform random sample of the training set, preserving the global class
distribution at every client. This is the null condition: in
expectation, no client drift can accumulate, and FedProx is predicted
to provide no advantage.
\item \textbf{Dirichlet-$\alpha = 0.1$ (literature-standard non-IID).}
For each class, the proportion of its samples allocated to each of the
7 clients is drawn from a 7-dimensional Dirichlet distribution with
concentration parameter $\alpha = 0.1$, the now-canonical setting
introduced by \citet{hsu2019dirichlet}.
\item \textbf{Engineered \texttt{balanced\_paired\_7\_clients} (mechanism
experiment, this work).} A custom non-IID partition designed so that
every minority class is held by exactly two clients and the majority
class (\textit{melanocytic nevi}) is held by all seven. Composition is
given in Table~\ref{tab:partition}.
\item \textbf{Engineered \texttt{specialist\_7\_clients} (pairing-lever
check).} A sister partition in which every minority class is held by
exactly \emph{one} client; serves as a pre-registered isolation of the
pairing mechanism versus a generic label-skew effect.
\end{itemize}

% INSERT_FROM_ORIGINAL: lines 1320-1352 (the partition table
% \texttt{tab:partition} and \texttt{fig:partition-heatmap}). Plus
% lines 462-540 (Specialist-Partition Pre-Registration) — keep this as
% MPhil-strong methodological discipline.

\paragraph{System heterogeneity (computational-work imbalance).}

System heterogeneity in this dissertation is operationalised as
per-client variable local-epoch budgets, following the convention of
\citet{li2020fedprox} §5.2; we do not simulate hardware latency or
network conditions. Three conditions are tested in the Flower runtime:

\begin{itemize}
\item \textbf{C0 (uniform compute):} every client performs $E = 20$
local epochs per round. This is the runtime-matched baseline against
which C1 and C2 are compared and is run separately from the
pure-PyTorch headline because cross-runtime mixing in inferential
contrasts is avoided (§\ref{sec:baseline-fl}).
\item \textbf{C1 (fixed stragglers):} clients $C_5$ and $C_6$ perform
only $E_{\text{straggler}} = 5$ epochs per round; the remaining five
clients perform $E = 20$. This is descriptive only — the fixed
identity of $C_5$ partially removes the melanoma drift that the
headline analysis identifies as the per-class FedProx driver, so this
condition is reported but not used for primary inference.
\item \textbf{C2 (random stragglers):} 4 of 7 clients per round are
randomly designated stragglers with
$E_i \sim \mathrm{Uniform}\{1, 2, \ldots, 19\}$. This is the
primary system-heterogeneity condition for inference.
\end{itemize}

%==============================================================================
\subsection{FedProx}
\label{sec:fedprox}
%==============================================================================

% INSERT_FROM_ORIGINAL: condensed version of lines 1080-1268 focused
% on FedProx (drop the FedNova content; that gets its own subsection
% below). Aim for ~1 page.

To mitigate the challenges introduced by statistical and system
heterogeneity, FedProx \citep{li2020fedprox} extends the standard
FedAvg approach by adding a proximal term to the local objective on
each client, penalising deviation from the current global model.
In FedProx, each client $k$ minimises the following local objective:

\begin{equation}
\min_w \;\; F_k(w) \;+\; \frac{\mu}{2}\,\bigl\lVert w - w^{(t)} \bigr\rVert^2,
\label{eq:fedprox-method}
\end{equation}

where $F_k(w)$ is the local cross-entropy loss for client $k$, $w$ are
the local model parameters, $w^{(t)}$ are the broadcast global model
parameters at round $t$, and $\mu \geq 0$ controls the strength of the
regularisation. At $\mu = 0$, FedProx reduces exactly to FedAvg. After
local training, the server aggregates the returned parameters using
the same size-weighted average as FedAvg (Eq.~\ref{eq:fedavg-method}).

We use $\mu = 0.01$, selected on a held-out three-seed pilot sweep
over $\{0.001, 0.01, 0.1\}$ on a different, discarded partition
design; $\mu$ was \emph{not} tuned on the test set of the final
10-seed sweep nor on the \texttt{balanced\_paired\_7\_clients}
partition. This protects the headline test comparison from
hyperparameter selection bias.

% INSERT_FROM_ORIGINAL: optionally the FedProx algorithm pseudocode
% block from lines 1123-1155 if you want to retain the formal
% presentation.

%==============================================================================
\subsection{FedNova}
\label{sec:fednova}
%==============================================================================

% INSERT_FROM_ORIGINAL: condensed version of FedNova content from
% lines 1080-1268 + lines 2192-2230 (which provides additional FedNova
% technical detail).

FedNova \citep{wang2020fednova} addresses a complementary axis of
heterogeneity to FedProx: when clients differ in the amount of local
computation they perform per round (heterogeneous $\tau_i$, here
driven by variable local epochs under C1 and C2), the naïve
size-weighted aggregation rule is \emph{objective-inconsistent}, and
FedNova removes the bias by normalising each client's update by its
effective number of local SGD steps.

For local SGD with momentum coefficient $m$ and $\tau_i$ local steps
at client $i$, the FedNova normaliser is

\begin{equation}
a_i \;=\; \sum_{j=1}^{\tau_i} \frac{1 - m^j}{1 - m}
       \;=\; \frac{\tau_i (1 - m) - m\,(1 - m^{\tau_i})}{(1-m)^2},
\label{eq:fednova-ai-method}
\end{equation}

which reduces to $a_i = \tau_i$ for vanilla SGD ($m = 0$). The
aggregation step is $d_i = w^t - w_i^{t+1}$,
$d_\text{norm} = \sum_i p_i\,d_i / a_i$,
$a_\text{eff} = \sum_i p_i a_i$, and
$w^{t+1} = w^t - a_\text{eff}\,d_\text{norm}$. With $m = 0$ \emph{and}
uniform $\tau_i$, FedNova reduces to FedAvg. Under $m = 0.9$ (the
local optimiser used throughout this work) even uniform $\tau_i$
produces a different aggregate from FedAvg, so the runtime-matched C0
baseline is run separately for FedNova rather than borrowed from the
FedAvg run.

The implementation is brute-force verified against canonical reference
values from \citet{wang2020fednova} (e.g.\ for $\tau = 10$, $m = 0.9$:
$a_i = 41.381$); see §\ref{sec:protocol} for the unit-test inventory.

%==============================================================================
\chapter{Results and Discussion}
\label{ch:results}
%==============================================================================

This chapter presents the experimental findings, starting with the
centralised baseline and progressing through federated-learning
scenarios affected by statistical and system heterogeneity. Particular
attention is given to evaluating the performance of FedProx and
FedNova across these settings. Methodological observations and
mechanism-interpretation discussion are interleaved with the
empirical findings rather than relegated to a separate Discussion
chapter.

%==============================================================================
\section{Centralised baseline}
\label{sec:centralised-results}
%==============================================================================

% INSERT_FROM_ORIGINAL: lines 1865-1885 (the "Centralised upper bound
% (3-seed sweep)" paragraph) + lines 1938-1965 (the federation-tax
% prose). Drop forward references to FL results.

The centralised baseline anchors the upper bound on test macro-F1
achievable when all training data is pooled into a single non-federated
run. Using the \texttt{DermMNISTCNN} architecture described in
§\ref{sec:centralised-method}, the centralised model reaches test
macro-F1 = $\mathbf{0.562 \pm 0.006}$ across three seeds
($\{42, 123, 456\}$); per-seed values 0.555, 0.568, 0.562. The
across-seed standard deviation (0.006) is markedly smaller than the
federated across-seed standard deviations reported in
§\ref{sec:engineered-fl} (FedAvg 0.025, FedProx 0.014), as expected:
centralised training has no client-aggregation noise.

% INSERT_FROM_ORIGINAL: lines 1938-1965 (the federation-tax discussion
% — relative gap recovery, per-class concentration on melanoma + BCC,
% \textit{melanocytic nevi} reaching centralised-level under both
% federated algorithms).

%==============================================================================
\section{Federated Learning}
\label{sec:fl-results}
%==============================================================================

%------------------------------------------------------------------------------
\subsection{Initial Setup: Engineered Paired Partition}
\label{sec:engineered-fl}
%------------------------------------------------------------------------------

% Headline result on the engineered 7-client partition.

\paragraph{Analysis of the global model.}

% INSERT_FROM_ORIGINAL: lines 1414-1497 (Final Test Performance: the
% headline paired Wilcoxon, Hodges-Lehmann, LOSO, per-seed table
% \texttt{tab:per-seed}, headline table \texttt{tab:headline}).

% INSERT_FROM_ORIGINAL: lines 1596-1653 (Convergence Behaviour:
% \texttt{fig:curves-main} and the global-model trajectories) +
% lines 1701-1732 (Summary of Numerical Observations from the
% Curves: \texttt{tab:curve-stats}).

% INSERT_FROM_ORIGINAL: lines 1779-1805 (Peak vs Final Round:
% \texttt{tab:best-vs-last}).

\paragraph{Analysis of individual clients and classes.}

% INSERT_FROM_ORIGINAL: lines 1498-1594 (Per-Class Test Performance:
% per-class table \texttt{tab:per-class}, worst-case
% \texttt{tab:worst-per-class}, and the Holm-corrected discussion).

% INSERT_FROM_ORIGINAL: lines 1655-1699 (Per-Class Convergence:
% \texttt{fig:curves-per-class} and \texttt{fig:per-client-curves}).

% INSERT_FROM_ORIGINAL: lines 1734-1777 (Per-Specialty Analysis:
% \texttt{tab:per-specialty}, the three specialty pairs + generalist
% asymmetry).

\paragraph{Mechanism interpretation.}

% Interleave the Discussion content from lines 2041-2097 here.
% Specifically the three paragraphs:
%   (i) training-objective panel caveat
%   (ii) per-class / per-specialty mechanism asymmetry
%   (iii) mid-training communication-efficiency gain + peak-vs-final
%        plateau evidence
% INSERT_FROM_ORIGINAL: lines 2041-2097.

%------------------------------------------------------------------------------
\subsection{Statistical Heterogeneity (Multi-Partition Comparison)}
\label{sec:stat-het-multi}
%------------------------------------------------------------------------------

% Multi-partition sweep (IID, Dirichlet, specialist, paired). Lead
% with the central comparison figure/table, then per-condition prose.

% INSERT_FROM_ORIGINAL: lines 1807-1864 (the partition-robustness
% framing prose, the three bullet points IID/Dirichlet/paired, +
% \texttt{fig:federation-tax}).

% INSERT_FROM_ORIGINAL: lines 1887-1899 (IID falsification check prose
% with \TODOhpc placeholders for the actual numbers).

% INSERT_FROM_ORIGINAL: lines 1901-1915 (Dirichlet-$\alpha = 0.1$ prose
% with \TODOhpc placeholders).

% INSERT_FROM_ORIGINAL: lines 1918-1936 (the
% \texttt{tab:robustness} table consolidating all four conditions).

% NEW_CONTENT_REQUIRED: short paragraph summarising the four-point
% dose-response (IID → Dirichlet → specialist → paired) once all the
% sweeps land. Reference §05_when_fedprox_helps.tex (the synthesis
% subsection you drafted earlier).

%------------------------------------------------------------------------------
\subsection{System Heterogeneity}
\label{sec:sys-het}
%------------------------------------------------------------------------------

% The 7-client C0/C1/C2 sweep + FedNova arm. The data for this
% section is currently running on RunPod; populate when complete.

% NEW_CONTENT_REQUIRED: the system-heterogeneity results section.
% Structure:
%   (a) C0 baseline (Flower runtime-matched) — describe FedAvg / FedProx /
%       FedNova on uniform $E = 20$, all 7 clients full participation.
%       Reference the bug-fix in commit 40c91f4 (numpy seed overflow at
%       seed 8675309) and confirm the data was clean after re-run.
%   (b) C1 fixed-stragglers — descriptive only, note the C5/C6
%       confound described in §2.4.2.
%   (c) C2 random-stragglers — primary inferential condition; the H2
%       between-condition contrast (Δ_C2 vs Δ_C0) is tested by paired
%       Wilcoxon under Bonferroni α/2 = 0.025.
%   (d) FedNova in C2 — does the objective-inconsistency correction
%       help on top of FedProx's proximal anchor?
%   (e) Mechanism interpretation: when does FedProx's proximal anchor
%       help vs. when does FedNova's normalised aggregator help.

System heterogeneity results pending; sweeps currently executing on
RunPod (4090) and CSD3 (A100). \TODOhpc{full results}

%------------------------------------------------------------------------------
\subsection{Summary and Comparison of Heterogeneity Experiments}
\label{sec:summary-het}
%------------------------------------------------------------------------------

% INSERT_FROM_ORIGINAL: lines 1967-2037 (the summary section + the
% across-conditions \texttt{tab:summary} table).

%==============================================================================
\section{Hardware and Resource Limitations}
\label{sec:hardware}
%==============================================================================

% INSERT_FROM_ORIGINAL: lines 2332-2427 (the full Hardware section
% verbatim, including \texttt{tab:resources}).

%==============================================================================
\chapter{Conclusion}
\label{ch:conclusion}
%==============================================================================

% Fold Future Work INTO the Conclusion (~1.5 pages with future work as
% the closing paragraphs). Structure:
%
%  - Opening paragraph: restate the clinical-and-methodological
%    motivation (cancer detection + privacy + FL).
%  - Middle paragraphs: restate the calibrated thesis-level claim
%    (FedProx improves macro-F1 on the engineered partition; melanoma
%    is the only Holm-significant per-class gain; the result transfers
%    to Dirichlet-α=0.1 with attenuated magnitude; FedProx is null on
%    IID as predicted by mechanism; FedNova provides complementary
%    benefit on C2 system heterogeneity).
%  - Acknowledgement of scope: simulated label-skew not real
%    multi-site; 28×28 not clinical resolution; full participation
%    only; 10 paired seeds.
%  - Future-work paragraph (integrated rather than separate chapter):
%    64×64 resolution sensitivity; partial participation; FedLC/
%    class-weighted CE/focal loss baselines; SCAFFOLD; update-norm
%    logging; real multi-site HAM10000 source-id partition.
%  - Closing sentence reinforcing the main result.

% INSERT_FROM_ORIGINAL: lines 2099-2115 (the "Calibrated thesis-level
% claim" block quote) — this is the cleanest single statement of your
% thesis claim and belongs in the Conclusion.

\TODOhpc{Write the Conclusion (~500-700 words) once the system-het
sweep lands. Use the calibrated-claim block-quote from lines
2099-2115 of the original draft as the spine; expand into 4-5
paragraphs.}

%==============================================================================
% References
%==============================================================================
\bibliographystyle{plainnat}
\bibliography{references}

\end{document}
